\documentclass[a4paper,12pt]{article}
 
% ------------------- Encoding & Fonts --------------------------------------
\usepackage[utf8]{inputenc}
\usepackage[T1]{fontenc}
\usepackage[ngerman,provide=*]{babel}
\usepackage{microtype}          % Bessere Typografie
\microtypesetup{expansion=false} % Deaktiviert Fontexpansion (CM-Fonts)
 
% ------------------- Seitenlayout ------------------------------------------
\usepackage{geometry}
\geometry{a4paper, margin=25mm, top=30mm, bottom=28mm, headsep=8mm, headheight=15pt}
 
% ------------------- Kopf- und Fußzeilen -----------------------------------
\usepackage{fancyhdr}
\pagestyle{fancy}
\fancyhf{}
\fancyhead[L]{\small\itshape Campus Next-Gen Data-Hub~--~Zwischenbericht}
\fancyhead[R]{\small\itshape Gruppe Airbyte}
\fancyfoot[C]{\small\thepage}
\renewcommand{\headrulewidth}{0.4pt}
 
% ------------------- Farben ------------------------------------------------
\usepackage{xcolor}
\definecolor{hsoBlue}{RGB}{0,83,143}
\definecolor{tableHead}{RGB}{220,232,246}
\definecolor{statusGreen}{RGB}{34,139,34}
\definecolor{statusOrange}{RGB}{200,100,0}
\definecolor{codeBox}{RGB}{248,248,250}
 
% ------------------- Tabellen ----------------------------------------------
\usepackage{booktabs}
\usepackage{array}
\usepackage{colortbl}
\usepackage{float}
 
\newcolumntype{L}[1]{>{\raggedright\let\newline\\\arraybackslash}p{#1}}
\newcolumntype{R}[1]{>{\raggedleft\let\newline\\\arraybackslash}p{#1}}
\newcolumntype{C}[1]{>{\centering\let\newline\\\arraybackslash}p{#1}}
 
% ------------------- Captions ----------------------------------------------
\usepackage{caption}
\captionsetup{font=small, labelfont=bf, skip=5pt}
\captionsetup[table]{position=top}
 
% ------------------- Listen ------------------------------------------------
\usepackage{enumitem}
 
% ------------------- Verbatim/Code -----------------------------------------
\usepackage{fancyvrb}
 
% ------------------- Grafik ------------------------------------------------
\usepackage{graphicx}
 
% ------------------- Links -------------------------------------------------
\usepackage{hyperref}
\hypersetup{
    colorlinks  = true,
    linkcolor   = hsoBlue,
    urlcolor    = hsoBlue,
    citecolor   = black,
    breaklinks  = true,
    pdftitle    = {Zwischenbericht -- Campus Next-Gen Data-Hub},
    pdfauthor   = {Lahres Timo, Braeutigam Rebecca, Horst Isabella},
    pdfsubject  = {Evaluation von Airbyte als ETL-Werkzeug}
}
 
% ------------------- Symbole -----------------------------------------------
\usepackage{pifont}
\newcommand{\done}{\textcolor{statusGreen}{\ding{52}}}      % Haken (erledigt)
\newcommand{\inprog}{\textcolor{statusOrange}{\ding{119}}}  % Kreis (laufend)
 
% ------------------- Absatzformatierung ------------------------------------
\setlength{\parskip}{0.6em}
\setlength{\parindent}{0em}
\setlength{\emergencystretch}{2em}  % Verhindert kleine Overfull-Hboxes
 
% ------------------- Metadaten ---------------------------------------------
\newcommand{\reportdate}{07.06.2026}
 
% ===========================================================================
\begin{document}
% ===========================================================================
 
% --- Titelseite -------------------------------------------------------------
\begin{titlepage}
    \centering
    \vspace*{3.5cm}
    {\LARGE\bfseries Zwischenbericht\par}
    \vspace{1.2em}
    {\color{hsoBlue}\rule{\linewidth}{1.5pt}}\par
    \vspace{1em}
    {\Large\bfseries Campus Next-Gen Data-Hub\par}
    \vspace{0.6em}
    {\large Evaluation von Airbyte als ETL- und Integrationswerkzeug\par}
    \vspace{0.4em}
    {\large Ablösung von Talend\par}
    \vspace{1em}
    {\color{hsoBlue}\rule{\linewidth}{0.5pt}}\par
    \vspace{3em}
    {\large INF-M Modul Projekte SoSe~'26\par}
    \vspace{0.5em}
    {\large Gruppe Airbyte\par}
    \vspace{2.5em}
    {\large
        Lahres Timo\\[0.5em]
        Bräutigam Rebecca\\[0.5em]
        Horst Isabella
    \par}
    \vspace{3em}
    {\large \reportdate\par}
    \vfill
\end{titlepage}
 
% --- Inhaltsverzeichnis ------------------------------------------------------
\thispagestyle{empty}
\tableofcontents
\clearpage
 
% ===========================================================================
\section{Projektziel \& Kontext}
% ===========================================================================
 
Im Projekt evaluieren wir \href{https://airbyte.com/}{Airbyte} als modernes ETL- und
Datenintegrationswerkzeug mit dem Ziel, das bisher eingesetzte \textbf{Talend} in der
Hochschul-IT abzulösen.
 
Die Evaluationsumgebung läuft \textbf{lokal in Docker Desktop} und stellt einen
realistischen Ausschnitt der Hochschul-Datenlandschaft dar, einschließlich anonymisierter
Studierenden-, Gebäude-, Instituts- und Personaldaten.
 
Die Evaluation ist in \textbf{sechs Testszenarien} gegliedert:
\begin{itemize}[noitemsep, leftmargin=*]
    \item DB-Connector
    \item Facility-Management-Sync
    \item Bild-BLOBs
    \item Studenten/Personal-Mapping
    \item Incremental Sync / IdM
    \item Web-APIs
\end{itemize}
 
Alle Szenarien sind detailliert in
\href{https://github.com/Timbo3399/INFM_Airbyte/blob/main/docs/testszenarien.md}{testszenarien.md}
beschrieben.
 
% ===========================================================================
\section{Erreichte Meilensteine (Stand 06.06.2026)}
% ===========================================================================
 
\begin{table}[H]
    \caption{Meilenstein-Übersicht mit Bearbeitungsstand}
    \centering\small
    \begin{tabular}{@{}L{6.2cm}C{3.2cm}L{5cm}@{}}
        \toprule
        \rowcolor{tableHead}
        \textbf{Meilenstein} & \textbf{Status} & \textbf{Beleg / Doku} \\
        \midrule
        Installation des Systems &
            \done~DB-Stack + Airbyte lauffähig &
            \href{https://github.com/Timbo3399/INFM_Airbyte/blob/main/docs/installation-guide.md}{installation-guide.md} \\
        \addlinespace[4pt]
        Zugang für Betreuer &
            \inprog~dokumentiert, Einrichtung im Termin &
            \href{https://github.com/Timbo3399/INFM_Airbyte/blob/main/docs/zugang.md}{zugang.md} \\
        \addlinespace[4pt]
        Einfacher ETL-Prozess &
            \done~durchgeführt und verifiziert (Postgres\,$\rightarrow$\,Postgres) &
            \href{https://github.com/Timbo3399/INFM_Airbyte/blob/main/docs/etl-prozess.md}{etl-prozess.md} \\
        \addlinespace[4pt]
        Beginn der Dokumentation &
            \done~README + 8 Dokumente unter \texttt{docs/} &
            dieses Repo \\
        \bottomrule
    \end{tabular}
\end{table}
 
% --- Installation ------------------------------------------------------------
\subsection{Installation}
 
Das Setup ist vollständig skriptbasiert und reproduzierbar:
\begin{itemize}[noitemsep, leftmargin=*]
    \item \texttt{scripts/install.ps1} (Windows) bzw.\ \texttt{scripts/install.sh}
          (Linux/macOS) startet den kompletten Datenbank-Stack (Source-PostgreSQL,
          Ziel-PostgreSQL, Ziel-MySQL, CSV-File-Server), wartet auf den
          \texttt{healthy}-Status und lädt die Testdaten.
    \item \texttt{scripts/setup-airbyte.ps1} / \texttt{scripts/setup-airbyte.sh}
          installiert Airbyte Community Edition über das offizielle CLI \texttt{abctl}
          in einem lokalen Kubernetes-Cluster (\texttt{kind}) innerhalb von Docker Desktop.
    \item \textbf{Plattformen:} Das Setup steht für \textbf{Windows, Linux und macOS} bereit
          (identische Logik, plattformspezifische Skripte).
\end{itemize}
Alle vier Datenbank-/Server-Container laufen verifiziert im Zustand \texttt{healthy}.
 
% --- Datenbasis --------------------------------------------------------------
\subsection{Datenbasis (Source-PostgreSQL)}
 
Die anonymisierten Testdaten sind geladen:
 
\begin{table}[H]
    \caption{Tabellen in der Source-Datenbank und ihre Lademechanismen}
    \centering\small
    \begin{tabular}{@{}L{4.8cm}R{2.2cm}L{6.5cm}@{}}
        \toprule
        \rowcolor{tableHead}
        \textbf{Tabelle} & \textbf{Zeilen} & \textbf{Lademechanismus} \\
        \midrule
        \texttt{fm\_gebaeude}  &      25 & \texttt{scripts/load\_fm\_gebaeude.py} \\
        \texttt{fm\_inst}      &   2.083 & \texttt{scripts/load\_fm\_inst.py} \\
        \texttt{k\_plz}        &  34.172 & \texttt{scripts/load\_k\_plz.py} \\
        \texttt{fm\_rna}       &     379 & \texttt{scripts/load\_json.py} \\
        \texttt{hso\_personal} &     870 & \texttt{scripts/load\_json.py} \\
        \texttt{hso\_students} &       0 &
            nur via File-Connector (s.\ Abschn.~\ref{sec:file-connector}) \\
        \texttt{fm\_stamm}     &       0 &
            keine Quelldatei (s.\ Abschn.~\ref{sec:fm-stamm}) \\
        \bottomrule
    \end{tabular}
\end{table}
 
% --- Connectoren & ETL -------------------------------------------------------
\subsection{Angelegte Airbyte-Connectoren \& erster ETL-Lauf}
\label{sec:file-connector}
 
In Airbyte sind folgende Connectoren angelegt und per Verbindungstest erfolgreich geprüft:
 
\begin{itemize}[noitemsep, leftmargin=*]
    \item \textbf{Sources (5):}
    \begin{itemize}[noitemsep, leftmargin=*]
        \item \texttt{HSO Source PostgreSQL}
              (Postgres, Update-Methode \textit{User Defined Cursor})
        \item \texttt{HSO CSV hso\_students} (\texttt{local}, \texttt{/local/*.csv})
        \item \texttt{HSO CSV k\_plz} (\texttt{local}, \texttt{/local/*.csv})
        \item \texttt{HSO CSV fm\_gebaeude} (\texttt{local}, \texttt{/local/*.csv})
        \item \texttt{HSO CSV fm\_inst} (\texttt{local}, \texttt{/local/*.csv})
    \end{itemize}
    \item \textbf{Destinations (2):}
    \begin{itemize}[noitemsep, leftmargin=*]
        \item \texttt{HSO Dest PostgreSQL} (Port 5434)
        \item \texttt{HSO Dest MySQL}
              (Port 3306, SSL aus, \texttt{allowPublicKeyRetrieval=true},
              Raw-DB \texttt{destdb})
    \end{itemize}
\end{itemize}
 
Es wurden \textbf{drei Connections} (jeweils \textit{Full Refresh\,|\,Overwrite})
ausgeführt und das Ergebnis \textbf{direkt in der jeweiligen Ziel-DB} verifiziert:
 
\begin{table}[H]
    \caption{Ausgeführte ETL-Connections und verifizierte Ergebnisse}
    \centering\small
    \begin{tabular}{@{}L{6.8cm}C{2.8cm}C{3.1cm}@{}}
        \toprule
        \rowcolor{tableHead}
        \textbf{Connection} & \textbf{Streams} & \textbf{Ergebnis (Ziel-DB)} \\
        \midrule
        \texttt{HSO Source PostgreSQL\,$\rightarrow$\,HSO Dest PostgreSQL} &
            \texttt{fm\_gebaeude},\newline\texttt{k\_plz} &
            25\,/\,34.172~\done \\
        \addlinespace[4pt]
        \texttt{HSO Source PostgreSQL\,$\rightarrow$\,HSO Dest MySQL} &
            \texttt{fm\_gebaeude},\newline\texttt{k\_plz} &
            25\,/\,34.172~\done \\
        \addlinespace[4pt]
        \texttt{HSO CSV hso\_students\,$\rightarrow$\,HSO Dest PostgreSQL} &
            \texttt{hso\_students}\newline(File) &
            \textbf{5.052 Zeilen}~\done \\
        \bottomrule
    \end{tabular}
\end{table}
 
\textbf{Bemerkenswert:} Der File-Connector lud die strukturell defekte
\texttt{hso\_students.csv} vollständig (5.052~Zeilen) in die DB~-- dieselbe Datei, an
der ein direktes PostgreSQL-\texttt{COPY} scheiterte (0~Zeilen,
s.\ Abschn.~\ref{sec:datenqualitaet}). Pandas im File-Connector toleriert die
Spalteninkonsistenzen. Damit sind \textbf{DB-Connector} (PG\,$\rightarrow$\,PG,
PG\,$\rightarrow$\,MySQL) und \textbf{File-Connector} (CSV\,$\rightarrow$\,DB)~-- der
Kern von Szenario~1 inkl.\ \glqq{}PostgreSQL\,$\rightarrow$\,MySQL\grqq{}~-- nachgewiesen.
 
Zusätzlich stellt der Dienst \textbf{PostgREST} (\texttt{hso\_postgrest}, Szenario~6a)
eine REST-API auf die Ziel-DB bereit. Der Aufruf
\texttt{GET http://localhost:3000/k\_plz?limit=5} liefert die synchronisierten Daten als JSON.
 
% ===========================================================================
\section{Architektur (Kurzüberblick)}
% ===========================================================================
 
Die vollständige Beschreibung findet sich in
\href{https://github.com/Timbo3399/INFM_Airbyte/blob/main/docs/architektur.md}{architektur.md}.
Alle Dienste laufen in Docker Desktop in einem gemeinsamen Netzwerk
(\texttt{airbyte\_net}); Airbyte selbst läuft in einem \texttt{kind}-Kubernetes-Cluster
und erreicht die Datenbanken über \texttt{host.docker.internal}.
 
\newsavebox{\archbox}
\begin{lrbox}{\archbox}
\begin{BVerbatim}[fontsize=\small, baselinestretch=1.15]
+--------------------------------------------------------------+
|                     Docker Desktop                           |
|                                                              |
|  +-----------------+   +------------------------------+    |
|  | source-postgres  |   |   Airbyte (kind-Cluster)     |    |
|  | Testdaten        +-->|                              |    |
|  | localhost:5433   |   |   UI: localhost:8000         |    |
|  +-----------------+   |   Connector-Pods             |    |
|                         |                              |    |
|  +-----------------+   |                              |    |
|  | file-server      +-->|                              |    |
|  | localhost:8888   |   +-----------+------------------+    |
|  +-----------------+               |                       |
|                                    | (schreibt)             |
|  +-----------------+               |                       |
|  | dest-postgres    |<--------------+                       |
|  | localhost:5434   |               |                       |
|  +-----------------+               |                       |
|  +-----------------+               |                       |
|  | dest-mysql       |<--------------+                       |
|  | localhost:3306   |                                      |
|  +-----------------+                                      |
+--------------------------------------------------------------+
\end{BVerbatim}
\end{lrbox}
 
\begin{figure}[H]
    \centering
    \setlength{\fboxsep}{8pt}%
    \colorbox{codeBox}{\usebox{\archbox}}
    \caption{Systemarchitektur: Docker-Desktop-Netzwerk mit Airbyte im kind-Cluster}
\end{figure}
 
% ===========================================================================
\section{Besonderheit: Datenqualität der Quell-CSVs}
\label{sec:datenqualitaet}
% ===========================================================================
 
Die bereitgestellten CSV-Dateien ließen sich \textbf{nicht} per direktem
PostgreSQL-\texttt{COPY} laden (Details in Abschn.~\ref{sec:loesungen-datenqualitaet}).
Wir haben daher tolerante Python-Loader implementiert, die die Daten nach dem
Containerstart bereinigt einspielen. Die Datei \texttt{hso\_students.csv} ist
strukturell so inkonsistent, dass sie aktuell nicht zuverlässig in die relationale
Source-DB geladen werden kann; Studierendendaten werden stattdessen über den
\textbf{Airbyte File-Connector} als Flatfile-Quelle eingebunden.
 
% ===========================================================================
\section{Probleme \& Lösungen}
% ===========================================================================
 
% --- Windows -----------------------------------------------------------------
\subsection{Windows-/Umgebungsspezifische Hürden}
 
\begin{table}[H]
    \caption{Windows-spezifische Hürden und ihre Lösungen}
    \centering\small
    \begin{tabular}{@{}L{5cm}L{4.5cm}L{5cm}@{}}
        \toprule
        \rowcolor{tableHead}
        \textbf{Problem} & \textbf{Ursache} & \textbf{Lösung} \\
        \midrule
        \texttt{install.ps1} meldete „Python gefunden", JSON-Laden schlug aber fehl &
            \texttt{python} ist unter Windows oft nur der Microsoft-Store-Platz\-halter;
            \texttt{Get-Command} findet ihn, er liefert aber keine echte Version &
            Echte Versionsprüfung (\texttt{py}/\texttt{python}/\texttt{python3});
            zusätzlich \textbf{Docker-Fallback}, der die Daten ganz ohne Host-Python lädt \\
        \addlinespace[4pt]
        \texttt{setup-airbyte.ps1} fand kein abctl-Asset &
            Falsches Namensschema~-- korrekt ist
            \texttt{abctl-<version>-}\linebreak\texttt{windows-<arch>.zip};
            zudem liegt \texttt{abctl.exe} in einem Unterordner &
            Asset per Muster ermitteln, aus Unterordner entpacken;
            Architektur-Erkennung PowerShell-7-fest \\
        \addlinespace[4pt]
        File-Server-Container dauerhaft \texttt{unhealthy}, Setup lief in Timeout &
            Healthcheck nutzte \texttt{localhost} $\to$ im Container zuerst IPv6
            (\texttt{::1}), nginx lauscht aber nur auf IPv4 &
            Healthcheck auf \texttt{127.0.0.1} umgestellt \\
        \bottomrule
    \end{tabular}
\end{table}
 
% --- Datenqualität -----------------------------------------------------------
\subsection{Datenqualität der Quell-CSVs (Hauptproblem)}
\label{sec:loesungen-datenqualitaet}
 
Der SQL-Init lud die drei Kern-Tabellen \textbf{gar nicht}~-- \texttt{COPY} brach unter
\texttt{ON\_ERROR\_STOP} bereits an der ersten fehlerhaften Datei ab, wodurch alle
folgenden leer blieben. Die konkreten Probleme je Datei:
 
\begin{itemize}[noitemsep, leftmargin=*]
    \item \textbf{\texttt{fm\_gebaeude.csv}:} unquotierte Kommas in Textfeldern
          (z.\,B.\ „Hörsäle,Bib.,RZ"), eingebettete Header-Zeilen.
    \item \textbf{\texttt{k\_plz.csv}:} \textbf{3.417 Header-Zeilen} mitten in den
          Daten (zusammengesetzte Einzel-Exporte), Zeilen mit zu wenig Feldern,
          ein Feld (\texttt{krskfz}) breiter als das Schema.
    \item \textbf{\texttt{fm\_inst.csv}:} 86~Spalten (Tabelle nutzt 24),
          semikolon-getrennt, vereinzelt \textbf{NUL-Bytes},
          doppelt kodierte UTF-8-Umlaute.
    \item \textbf{\texttt{hso\_students.csv}:} \textbf{strukturell defekt}~--
          Datenzeilen haben mehr Spalten (ca.\ 50--56) als der eigene Header (40),
          dazu unbalanciertes Quoting und Float-formatierte Integer.
\end{itemize}
 
\textbf{Lösung:} Tolerante Python-Loader (\texttt{scripts/load\_*.py}) filtern Header,
reparieren bzw.\ überspringen fehlerhafte Zeilen, bereinigen NUL-Bytes und Mojibake und
führen Typkonvertierungen best-effort durch. Der SQL-\texttt{COPY}-Schritt wurde entfernt
(\texttt{01\_load\_data.sql}).
 
% --- Airbyte/abctl -----------------------------------------------------------
\subsection{Airbyte-/abctl-spezifische Hürden}
 
Airbyte Community läuft über \texttt{abctl} in einem \texttt{kind}-Kubernetes-Cluster.
Daraus ergaben sich mehrere nicht-triviale, in der offiziellen Dokumentation nicht
beschriebene Hürden:
 
\begin{table}[H]
    \caption{Airbyte-/abctl-spezifische Hürden und ihre Lösungen}
    \centering\small
    \begin{tabular}{@{}L{5cm}L{4.5cm}L{5cm}@{}}
        \toprule
        \rowcolor{tableHead}
        \textbf{Problem} & \textbf{Ursache} & \textbf{Lösung} \\
        \midrule
        \textbf{File-Connector} (\texttt{local}) findet \texttt{/local/*.csv} nicht &
            Connector-Pods (kind) sehen das Docker-Volume \texttt{oss\_local\_root} nicht &
            CSV-Verzeichnis beim Install via
            \texttt{abctl local install --volume "...:/local"} direkt in den
            Cluster mounten \\
        \addlinespace[4pt]
        \texttt{--volume} mit Windows-Pfad: \texttt{is not a valid volume spec} &
            \texttt{abctl} trennt den Volume-String stur an \texttt{:},
            der Laufwerks-Doppelpunkt (\texttt{C:}) kollidiert &
            Pfad in MSYS-Form \texttt{/c/Users/...} angeben \\
        \addlinespace[4pt]
        \textbf{Alle} Connector-Tests/Syncs hängen \texttt{Pending} &
            Aktiviertes lokales Volume erwartet PVC \texttt{airbyte-local-pvc}
            in jedem Job-Pod; es existierte nicht
            (\texttt{persistentvolumeclaim not found}) &
            PV (hostPath \texttt{/local}) + PVC \texttt{airbyte-local-pvc} anlegen;\newline
            \texttt{JOB\_KUBE\_LOCAL\_VOLUME\_}\linebreak\texttt{ENABLED=true} + Neustart
            von launcher/worker \\
        \addlinespace[4pt]
        \texttt{abctl local credentials} schlägt fehl
        (\texttt{invalid character '<'}) &
            Kombinierter Aufruf löst einen Org-Lookup aus,
            der HTML statt JSON liefert &
            E-Mail und Passwort in \textbf{zwei getrennten} Aufrufen setzen
            (erst \texttt{--email}, dann \texttt{--password}) \\
        \addlinespace[4pt]
        Connector-Auswahl heißt \textbf{„Postgres"} (nicht „PostgreSQL");
        Default-Update-Methode ist \textbf{CDC} &
            -- &
            In der UI „Postgres" wählen; Update-Methode auf
            \textit{User Defined Cursor} stellen
            (CDC bräuchte \texttt{wal\_level=logical}) \\
        \bottomrule
    \end{tabular}
\end{table}
 
Alle Lösungen sind in \texttt{scripts/setup-airbyte.ps1}/\texttt{.sh} automatisiert
und in
\href{https://github.com/Timbo3399/INFM_Airbyte/blob/main/docs/airbyte-setup.md}{airbyte-setup.md}
sowie
\href{https://github.com/Timbo3399/INFM_Airbyte/blob/main/docs/etl-prozess.md}{etl-prozess.md}
dokumentiert.
 
% ===========================================================================
\section{Offene Punkte / Fragen an die Betreuer}
\label{sec:offene-punkte}
% ===========================================================================
 
\begin{itemize}[leftmargin=*, itemsep=0.6em]
    \item \textbf{\texttt{hso\_students.csv}~-- Soll-Struktur:}
          Die Datei ist nicht eindeutig ladbar (mehr Datenspalten als Header-Einträge,
          defektes Quoting). Ist das Einlesen einer strukturell defekten Datei
          gewollt, oder wird eine korrekt lesbare Version bereitgestellt?
          Dies blockiert aktuell Szenario~4 (Account-Mapping in die Source-DB).
 
    \item \textbf{\texttt{fm\_stamm} (Raumstammdaten):}
          \label{sec:fm-stamm}
          Es liegt keine Quelldatei vor. Wird diese Tabelle benötigt, und woher
          sollen die Daten kommen?
 
    \item \textbf{Zugang für Betreuer:}
          Airbyte Community Edition ist im Wesentlichen Single-User und läuft auf
          \texttt{localhost}. Wie soll der Betreuer-Zugang erfolgen?
          Gemeinsamer Admin-Login während des Vor-Ort-Termins, oder ist ein
          Remote-Zugang (z.\,B.\ Tunnel/VM) gewünscht?
          (Vorschlag in
          \href{https://github.com/Timbo3399/INFM_Airbyte/blob/main/docs/zugang.md}{zugang.md}.)
 
    \item \textbf{Sync-Strategie:}
          Wir nutzen den Cursor-Modus (\texttt{updated\_at}) statt CDC, da CDC
          zusätzliche PostgreSQL-Konfiguration (\texttt{logical WAL},
          Replication~Slot) erfordert. Ist CDC für die Evaluation gewünscht?
 
    \item \textbf{Scope:}
          Welche der sechs Szenarien haben für die Bewertung Priorität?
\end{itemize}
 
% ===========================================================================
\section{Anforderungs- und Szenarien-Status}
% ===========================================================================
 
Eine vollständige Gegenüberstellung aller Kickoff-Anforderungen und der sechs
Szenarien mit Bewertung der Airbyte-Eignung und unserem Umsetzungsstand findet sich
in
\href{https://github.com/Timbo3399/INFM_Airbyte/blob/main/docs/anforderungen.md}{anforderungen.md}.
Die wesentlichen Befunde:
 
\begin{itemize}[noitemsep, leftmargin=*]
    \item \textbf{Erfüllt:} Open Source/Community,\linebreak PostgreSQL- \& MySQL-Anbindung,
          CSV/JSON/Excel-Dateien, Logging/Monitoring, einfacher ETL-Lauf (verifiziert).
    \item \textbf{Einschränkungen gegenüber Talend:}
    \begin{itemize}[noitemsep, leftmargin=*]
        \item Kein OSS-Connector für \textbf{Informix}
        \item Kein nativer \textbf{XML}-Connector
        \item Keine freie Code-Snippet-Ausführung
              (Mapping nur via dbt-SQL oder Custom-Connector)
    \end{itemize}
\end{itemize}
 
% ===========================================================================
\section{Nächste Schritte}
% ===========================================================================
 
\begin{itemize}[noitemsep, leftmargin=*]
    \item Szenario~1 abrunden: Sync nach MySQL + ein
          File\,$\rightarrow$\,DB-Sync (Nachweis File-Connector-Last).
    \item Weitere Szenarien umsetzen: FM-Denormalisierung (dbt/View),
          BLOB-Bilder, Incremental Sync (IdM), Web-APIs (PostgREST/SOAP).
    \item Klärung der offenen Fragen aus Abschn.~\ref{sec:offene-punkte}.
\end{itemize}
 
% ===========================================================================
\section{Anhang: Repository-Struktur}
% ===========================================================================
 
Siehe
\href{https://github.com/Timbo3399/INFM_Airbyte/blob/main/README.md}{README.md}.
Das vollständige Setup ist in unter 20 Minuten reproduzierbar:
\begin{itemize}[noitemsep, leftmargin=*]
    \item Windows: \texttt{scripts/install.ps1}~+~\texttt{scripts/setup-airbyte.ps1}
    \item Linux/macOS: \texttt{scripts/install.sh}~+~\texttt{scripts/setup-airbyte.sh}
\end{itemize}
 
\end{document}
 
